%==============================================================================
%  FINAL PROJECT PAPER  (IEEE Conference format) -- Category B
%  Profiling-Guided CUDA Optimization of 30 HeCBench Benchmarks via an AI Agent
%==============================================================================
\documentclass[conference]{IEEEtran}

\usepackage{cite}
\usepackage{amsmath,amssymb,amsfonts}
\usepackage{algorithmic}
\usepackage{graphicx}
\usepackage{textcomp}
\usepackage{xcolor}
\usepackage{listings}
\usepackage{url}
\usepackage{booktabs}
\usepackage{tabularx}
\usepackage{microtype}
\usepackage[hidelinks]{hyperref}
% Set PDF metadata explicitly so hyperref does not try to parse the
% multi-line \title (the \\ and \footnotesize otherwise trigger a
% "Missing number" error when hyperref auto-derives pdftitle).
\hypersetup{
  pdftitle={Profiling-Guided CUDA Optimization of 30 HeCBench Benchmarks with an AI Coding Agent},
  pdfauthor={Final Project, Parallel Programming Spring 2026}
}

\newcommand{\mainbodyendpage}{?}
\AtEndDocument{%
  \typeout{====================================================}%
  \typeout{[PAGE BUDGET] Main body + references end on page \mainbodyendpage\space (limit: 6, Appendix excluded).}%
  \typeout{====================================================}%
}

\lstset{
  basicstyle=\ttfamily\footnotesize,
  breaklines=true,
  frame=single,
  columns=fullflexible,
  showstringspaces=false,
  keepspaces=true
}

\newcolumntype{Y}{>{\raggedright\arraybackslash}X}

\begin{document}

\title{Profiling-Guided CUDA Optimization of 30 HeCBench
Benchmarks with an AI Coding Agent\\
{\footnotesize Final Project --- Parallel Programming, Spring 2026}}

% Authors: r14922119 王昊平 (Hao-Ping Wang), r14922031 曾柏瑜 (Po-Yu Tseng),
%          r14922000 吳善凱 (Shan-Kai Wu). Romanized below so the paper compiles
%          with pdfLaTeX; to print the Chinese names, load a CJK package
%          (e.g. xeCJK with XeLaTeX) and substitute them.
\author{%
\IEEEauthorblockN{Hao-Ping Wang, Po-Yu Tseng, Shan-Kai Wu}
\IEEEauthorblockA{r14922119, \enspace r14922031, \enspace r14922146 \\
National Taiwan University}%
}

\maketitle

%------------------------------------------------------------------------------
\begin{abstract}
We study how well an AI coding agent (Claude Code) can optimize real GPU kernels
when it is driven by a strict, profiling-first workflow rather than by free-form
``make it faster'' requests. Starting from 30 CUDA programs in the HeCBench
suite, we instructed the agent to (1) profile each kernel with Nsight Compute,
(2) classify its bottleneck (memory-bound, compute-bound, latency/occupancy-bound,
or atomic-contention), and only then (3) apply a bottleneck-matched optimization,
verifying numerical correctness against the original program before reporting any
speedup. All builds and timings use one NVIDIA V100 with CUDA~12.3, comparing the
agent's code against the unmodified baseline compiled by the \emph{same} compiler.
Of 30 benchmarks, the agent produced confirmed speedups on 15 (best cases:
xsbench 5.90$\times$, Floyd--Warshall 5.43$\times$, bilateral 2.9$\times$,
nbody 2.0$\times$), found no measurable improvement on 11, retained one regression
for analysis, and could not run 3 for lack of input data. The most instructive
results are the failures: we give a
profiling-grounded root-cause analysis of every case where an ``obvious''
optimization made the code slower, which is the central lesson of the project.
\end{abstract}

\begin{IEEEkeywords}
GPU computing, CUDA, performance profiling, AI coding agents, prompt engineering,
roofline analysis.
\end{IEEEkeywords}

\vspace{0.5em}
\noindent\fbox{\parbox{0.97\columnwidth}{\textbf{Project Category:}\quad
$\boxtimes$~\textbf{B} (AI-agent code optimization)\\[2pt]
\footnotesize AI-agent code optimization of an existing benchmark suite.}}

%------------------------------------------------------------------------------
\section{Introduction}
GPU kernels are notoriously easy to ``optimize'' in ways that make them slower:
raising occupancy can spill registers, parallelizing a serial loop can multiply
the total work, and fusing kernels can duplicate computation. A human expert
avoids these traps by first measuring \emph{why} a kernel is slow. Our question is
whether an AI coding agent can be made to follow the same discipline, and how it
behaves when an optimization that ``looks right'' is actually wrong.

We take 30 CUDA programs from HeCBench~\cite{hecbench} --- spanning data
reduction/scan, dense and sparse linear algebra, stencils, convolution, $N$-body,
and Monte-Carlo transport --- and drive an AI coding agent through a fixed
\emph{profile $\rightarrow$ classify $\rightarrow$ optimize $\rightarrow$ verify}
loop. The contribution is twofold: (1) a prompt-engineering design that forces the
agent to ground every change in profiler evidence and to verify correctness before
claiming a speedup; and (2) a candid root-cause analysis of the agent's failures,
which we argue is the most useful outcome of a Category-B study because it exposes
where naive ``parallelize harder'' reasoning breaks down.

%------------------------------------------------------------------------------
\section{Background and Related Work}
\textbf{Benchmark and dataset.}
HeCBench~\cite{hecbench} is a large open suite of heterogeneous-computing kernels.
We use 30 of its CUDA programs, chosen to cover distinct bottleneck
\emph{classes} rather than one application area: five reduction/compaction
programs, two linear solvers, seven stencils, five convolutions, six graph or
dynamic-programming workloads, and five simulation or Monte-Carlo programs.
The set includes both compact kernels such as scan and histogram and
application-scale workloads such as XSBench, miniWeather, and $N$-body.
We reuse the available input generators and built-in checks, adding output or
CPU-reference comparisons when a program has no self-check. Three programs
(\texttt{hotspot}, \texttt{hotspot3D}, \texttt{sobel}) require external input
files not present in the distribution (only DVC pointer stubs were available), so
they are reported as \emph{skipped}.

\textbf{Baseline and tooling.}
The baseline is each unmodified HeCBench program. The agent is Claude Code, given
shell access to a V100 cluster, the CUDA toolkit, Nsight Compute
(\texttt{ncu})/\texttt{nvprof}, and the benchmark sources. A key prior observation
is that an HeCBench reference run exists, but it was compiled with a different
toolchain; comparing against it would conflate compiler effects with the agent's
changes. We therefore always rebuild the baseline with the same compiler used for
the optimized code (Section~\ref{sec:setup}).

\textbf{Roofline intuition.}
The V100 has 80 SMs and $\approx$900\,GB/s of HBM2 bandwidth. A kernel is
\emph{memory-bound} if it is near that bandwidth, \emph{compute-bound} if SM
throughput is high while DRAM is idle, and \emph{latency/occupancy-bound} if both
are low and there are too few resident warps to hide latency. The optimization
that helps depends entirely on which of these holds --- the premise of our prompt.

%------------------------------------------------------------------------------
\section{Methodology}
The core of a Category-B project is how the agent is \emph{instructed}. Our prompt
design has four enforced stages; the agent is told not to advance until the
current stage is satisfied.

\textbf{Stage 1 --- Profile, do not guess.} For each benchmark the agent must run
\texttt{ncu} and extract SpeedOfLight (Compute\% vs.\ Memory\%), achieved vs.\
theoretical occupancy, waves/SM, L2 hit rate, and warp-stall reasons. ``Looks
slow'' is rejected as a reason; a metric must be cited.

\textbf{Stage 2 --- Classify the bottleneck.} The agent labels each kernel as
memory-, compute-, latency/occupancy-, or atomic-contention-bound, using one
rule: high Memory\%~$\Rightarrow$ widen loads / coalesce / cut redundant traffic;
high Compute\%~$\Rightarrow$ fewer instructions, \texttt{\_\_shfl} reductions,
fast intrinsics; both low~$\Rightarrow$ raise occupancy or cut launch/sync
overhead; high atomic traffic~$\Rightarrow$ warp-aggregate then one atomic.

\textbf{Stage 3 --- Optimize to match the bottleneck, and never increase total
work.} After several early regressions we added an explicit guardrail to the
prompt: \emph{prefer changes that remove waste (idle blocks, redundant loads,
host round-trips) over changes that add work to gain parallelism}. This single
instruction was the most important iteration on the agent (Appendix~D).

\textbf{Stage 4 --- Verify before claiming.} Every change must pass the program's
own check (or a CPU reference the agent writes) \emph{before} a speedup is
reported; a faster-but-wrong kernel is treated as a failure. The agent also
records \emph{rejected} experiments, not just successes.

We iterated the prompt three times: v1 (free-form ``optimize'') produced confident
but slower code; v2 added the profile-first requirement; v3 added the
no-extra-work guardrail and the correctness-gate-before-claim rule. All reported
numbers are from v3.

%------------------------------------------------------------------------------
\section{Experimental Setup}
\label{sec:setup}
\textbf{Hardware/software.} One NVIDIA Tesla V100-SXM2-32GB (Volta, sm\_70, 80 SM,
$\approx$900\,GB/s), CUDA~12.3, built \texttt{make ARCH=sm\_70}. Jobs run on a
dedicated GPU compute node via Slurm with a single pinned device
(\texttt{CUDA\_VISIBLE\_DEVICES=0}); running on the shared login node was found to
make timings drift by up to 2$\times$ and was abandoned.

\textbf{Metric.} Each program's own printed per-kernel/stage time (not wall time).
This is an average-per-launch figure, independent of the repeat count.

\textbf{Comparison.} Baseline and optimized are built with the \emph{same}
CUDA~12.3 compiler; speedup $=$ baseline\,/\,optimized. We use five-run medians
where repeated paired runs are available, and otherwise report the paired
measurement or measured range recorded by the benchmark batch. Some later
comparisons reuse five same-environment baseline logs because the printed metric
is already average-per-launch. Arguments are taken from what the baseline
actually ran (which is not always the Makefile \texttt{run:} target; e.g.\
\texttt{convolution3D} runs two real conv layers).

\textbf{Team organization.} The 30 programs were split among three members as
independent batches of 9, 5, and 16 benchmarks. Each member kept a per-benchmark
record of profile evidence, attempted changes, correctness, and timing; this paper
normalizes those records to the common criteria above.

\textbf{Correctness / success criterion.} A benchmark counts as a \emph{success}
only if it (i) passes the program's built-in verification (or a byte-identical
output/CPU-reference check for programs without one) \emph{and} (ii) achieves
$\geq$1.05$\times$ speedup. $0.95$--$1.05\times$ with correctness preserved is
``no measurable change''; $<0.95\times$ is a \emph{regression}; a wrong result is
a \emph{failure} regardless of speed.

%------------------------------------------------------------------------------
\section{Results}
Across 30 benchmarks, \textbf{15 achieved confirmed speedups}, \textbf{11 had no
measurable improvement}, \textbf{1 retained a regression for analysis}
(histogram), and \textbf{3 were skipped} for missing input data (hotspot,
hotspot3D, and sobel). These categories are mutually exclusive and account for
all 30 programs. Table~\ref{tab:all-results} gives one row per benchmark using
the metric recorded by its source report. A dash means that the source records a
speedup or profiling decision but not an absolute baseline/optimized value.

\begin{table*}[!t]
\caption{Complete 30-benchmark result summary. Speedup is baseline divided by
optimized time unless the row explicitly uses throughput (GFLOP/s).}
\label{tab:all-results}
\centering
\scriptsize
\setlength{\tabcolsep}{4pt}
\renewcommand{\arraystretch}{1.06}
\begin{tabularx}{\textwidth}{@{}r l r r r Y@{}}
\toprule
\textbf{\#} & \textbf{Benchmark / metric} & \textbf{Baseline} &
\textbf{Optimized} & \textbf{Speedup} & \textbf{Status / interpretation} \\
\midrule
1  & gaussian & 245.45\,ms & 142.70\,ms & \textbf{1.72$\times$} & Win \\
2  & thomas & 1.180\,ms & 0.975\,ms & \textbf{1.21$\times$} & Win \\
3  & jacobi & 40.2\,$\mu$s & 34.1\,$\mu$s & \textbf{1.18$\times$} & Win \\
4  & filter & 0.522\,ms & 0.528\,ms & 0.99$\times$ & Neutral; memory-bound \\
5  & histogram & 201.4\,$\mu$s & 388.0\,$\mu$s & \textbf{0.52$\times$} & Regression retained for analysis \\
6  & jaccard & 12.29\,ms & 12.30\,ms & 1.00$\times$ & Neutral; reverted to baseline algorithm \\
7  & bscan & 2894.7\,$\mu$s & 2909.0\,$\mu$s & 1.00$\times$ & Neutral \\
8  & atomicReduction & -- & -- & -- & Not modified; v4 reaches 863\,GB/s \\
9  & scan & 109586.96\,$\mu$s & 108237.26\,$\mu$s & 1.01$\times$ & Neutral \\
\midrule
10 & all-pairs-distance (k2) & 159.199\,$\mu$s & 87.916\,$\mu$s & \textbf{1.81$\times$} & Win; k3 is 1.82$\times$ \\
11 & nw & 21.923\,ms & 20.589\,ms & \textbf{1.06$\times$} & Win; five-run median \\
12 & snake (thresholds 0--25) & 6956.320\,$\mu$s & 3537.590\,$\mu$s & \textbf{1.97$\times$} & Win; aggregate metric \\
13 & bfs & 963.668\,$\mu$s & 962.199\,$\mu$s & 1.002$\times$ & Neutral; measurement noise \\
14 & floydwarshall & 9.815\,ms & 1.806\,ms & \textbf{5.43$\times$} & Win \\
\midrule
15 & convolution1D (int16) & 1752 & 1504 & \textbf{1.17--1.39$\times$} & Win; float/double are near 1.0$\times$ \\
16 & convolution3D & -- & -- & \textbf{1.33--1.69$\times$} & Win; range across heavy/small layers \\
17 & xsbench lookup & 0.342\,s & 0.058\,s & \textbf{5.90$\times$} & Win; lookup metric \\
18 & bilateral & -- & -- & \textbf{2.81--2.98$\times$} & Win; range across filter radii \\
19 & nbody (GFLOP/s) & 1957 & 3917 & \textbf{2.00$\times$} & Win; higher throughput is better \\
20 & stencil3d & -- & -- & $\sim$1.0$\times$ & Neutral; streaming coefficients \\
21 & convolutionSeparable & -- & -- & 1.006$\times$ & Neutral; measurement noise \\
22 & laplace3d & -- & -- & $\sim$1.0$\times$ & Neutral; partial-wave tail \\
23 & heat & -- & -- & $\sim$1.0$\times$ & Neutral; 88\% DRAM roofline \\
24 & lavaMD & -- & -- & \textbf{1.15$\times$} & Win \\
25 & fdtd3d & -- & -- & $\sim$1.0$\times$ & Neutral; problem too small \\
26 & miniWeather loop & 1.185\,s & 0.711\,s & \textbf{1.67$\times$} & Win \\
27 & srad compute stage & 0.0560\,s & 0.0530\,s & \textbf{1.06$\times$} & Win \\
28 & hotspot & -- & -- & -- & Skipped; input data unavailable \\
29 & hotspot3D & -- & -- & -- & Skipped; input data unavailable \\
30 & sobel & -- & -- & -- & Skipped; input image unavailable \\
\bottomrule
\end{tabularx}
\end{table*}

The wins cluster by \emph{mechanism}, not by application: the largest came from
algorithm/access-pattern changes (xsbench energy-sort, blocked Floyd--Warshall,
nbody tiling), the mid-range from removing structural waste (gaussian idle blocks,
miniWeather host round-trips, jacobi per-iteration syncs), and the small ones from
instruction-level cleanups (thomas CSE, convolution fast-paths). Every kernel
already at a memory or compute roofline yielded no honest speedup, while the
neutral BFS and Jaccard cases showed that reducing one local cost need not
improve the full kernel.

%------------------------------------------------------------------------------
\section{Discussion and Analysis}
The failures are the most informative part of the study. Every case where the
agent first made code \emph{slower} shares one root cause: an optimization that
added total work or overhead to chase a metric, instead of removing waste. The
profiler is what ultimately revealed each mistake.

\textbf{thomas (PCR attempt, 0.55$\times$).} The agent saw 9.95\% occupancy and
parallelized one tridiagonal system across a 1024-thread block using Parallel
Cyclic Reduction. PCR replaced a barrier-free, register-only serial sweep with
$\sim$20 block barriers and 48\,KB of shared memory per system. \emph{Root cause:}
occupancy was low but the SMs were not idle --- each thread was busy with a 1024-step
serial chain --- so adding synchronization cost more than the extra parallelism
returned. The agent reverted to the baseline algorithm and instead removed
redundant work (a recomputed denominator and repeated global loads), yielding a
safe 1.21$\times$.

\textbf{gaussian (fusion attempt, 0.85$\times$).} Fusing the pivot kernel into the
update kernel made every column thread recompute the same division
$a[\text{row}][t]/a[t][t]$, $\sim$4096$\times$ redundantly. \emph{Root cause:} the
saved kernel-launch overhead was far smaller than the duplicated arithmetic.
Re-aimed at the real waste --- a fixed full-size grid that launches blocks which
immediately return once the active submatrix shrinks --- shrinking the grid each
step gave 1.72$\times$ with the kernels left byte-identical.

\textbf{jaccard (binary-search attempt, 0.74$\times$).} To use the 7/8 idle lanes,
the agent replaced a two-pointer set-intersection ($O(N_i{+}N_j)$) with per-lane
binary search ($O(N_i\log N_j)$). \emph{Root cause:} spreading work over 8 lanes
does not pay for asymptotically \emph{more} total work; the baseline algorithm was
already efficient. Correctly parallelizing a two-pointer merge without adding work
is hard, so the agent reverted --- a legitimate ``no change'' outcome.

\textbf{histogram (grid-fill attempt, 0.52$\times$, retained).} The first kernel ran
at 20\% occupancy, so the agent enlarged the grid to fill the device. But more
blocks produce more partial histograms, and the \emph{second} (merge) kernel's
reduction loop grows with the number of partials. \emph{Root cause:} a two-stage
pipeline was optimized by looking only at stage one; the added merge cost exceeded
the occupancy gain. We retain this regression in the report as the clearest example
of a local metric (occupancy) misleading a global outcome.

\textbf{Rejected micro-optimizations.} A 128-bit vectorized
\texttt{convolution1D} regressed to 0.28$\times$ because the halo loads became
uncoalesced --- breaking coalescing cost far more than the redundant overlapped
reads it removed. A mirrored-atomic triangular variant of all-pairs k1 regressed
from atomic contention. Larger NW tiles and a fused/queue BFS frontier likewise
regressed. In every case the profiler explained why.

\textbf{Where the agent succeeded without prompting tricks.} The agent reliably
recognized roofline-bound kernels (Table~\ref{tab:all-results}) and chose
\emph{not} to change them --- an important behaviour, since a less disciplined
agent would have ``optimized'' a bandwidth-saturated stencil and reported noise as
a win. The profile-first prompt, plus the no-extra-work guardrail added in v3, were
what turned confident-but-wrong behaviour into measured, correctness-gated
results.

\textbf{Limitations.} Speedups are single-GPU and single-input. Repeated runs are
not available for every benchmark, so we do not report confidence intervals.
Three benchmarks could not be run for lack of data. The agent's correctness checks reuse
each program's own (sometimes weak) self-test; for programs without one we relied
on byte-identical output or a CPU reference the agent wrote.

%------------------------------------------------------------------------------
\section{Conclusion}
Driven by a strict profile-classify-optimize-verify prompt, an AI coding agent
optimized a 30-benchmark CUDA suite with results comparable to careful manual
tuning: 15 confirmed speedups (up to 5.9$\times$), 11 correct but neutral
outcomes, one retained regression, and three skipped inputs. The agent and the
profiler also caught and corrected several rejected regressions. The decisive
prompt-engineering lesson is a single rule: \emph{remove waste, do not add work};
without it the agent confidently produced slower code. The failures, analyzed by
root cause, show that the hard part of GPU optimization is not generating a
plausible change but knowing when a change that ``increases parallelism'' will
backfire --- a judgement that requires profiler evidence, which an agent can be
prompted to demand of itself.

%------------------------------------------------------------------------------
\begin{thebibliography}{00}
\bibitem{hecbench} Z.\ Jin and J.\ S.\ Vetter, ``HeCBench: A Benchmark Suite for
Heterogeneous Computing,'' Oak Ridge National Laboratory.
\url{https://github.com/ORNL/HeCBench}.
\bibitem{ncu} NVIDIA, ``Nsight Compute,'' CUDA Toolkit 12.3 documentation, 2023.
\bibitem{v100} NVIDIA, ``Tesla V100 GPU Architecture (Volta) Whitepaper,'' 2017.
\bibitem{harris} M.\ Harris and M.\ Garland, ``Optimizing Parallel Prefix
Operations for the Fermi Architecture,'' in \emph{GPU Computing Gems Jade
Edition}, 2012.
\end{thebibliography}

\xdef\mainbodyendpage{\thepage}

%==============================================================================
%  APPENDIX -- REQUIRED RAW MATERIALS (Category B)
%==============================================================================
\appendices

\section{Input Prompt (representative)}
\begin{lstlisting}
Data Compression / Reduction: histogram, scan, filter,
atomicReduction, bscan
Math: jacobi, gaussian, thomas, jaccard
These are the benchmarks I want to optimize; I want to beat
the benchmark. I need a workflow: bottleneck analysis ->
optimize. How to do bottleneck analysis -- use a profiler?

[follow-up, iterated over the session]
- pull the needed benchmarks from ORNL/HeCBench
- profile all first, then prioritize
- one benchmark at a time; record what each profile showed
  and what optimization was applied, then present results
- judge with the same compiler for baseline and optimized
  (do not let a newer nvcc give the baseline a free speedup)
\end{lstlisting}

\section{System Prompt (enforced workflow, paraphrased)}
\begin{lstlisting}
For each benchmark, do NOT skip steps:
1. PROFILE with ncu. Record Compute% vs Memory% (SpeedOfLight),
   achieved/theoretical occupancy, waves/SM, L2 hit, warp-stall
   reasons. "Looks slow" is not a reason; cite a metric.
2. CLASSIFY: memory- / compute- / latency-occupancy- / atomic-
   bound. Pick the matching lever:
     memory  -> wider/coalesced loads, cut redundant traffic
     compute -> fewer instr, __shfl reductions, fast intrinsics
     latency -> raise occupancy OR cut launch/sync overhead
     atomic  -> warp-aggregate (ballot+popc) then one atomic
3. OPTIMIZE, but PREFER REMOVING WASTE over ADDING WORK.
   Never increase total work to gain parallelism.
4. VERIFY against the program's own check (or a CPU reference)
   BEFORE reporting any speedup. A fast-but-wrong kernel is a
   failure. Record rejected experiments too.
Measure on a dedicated GPU node, single pinned device, same
CUDA version for baseline and optimized, with 5-run medians
where feasible and an explicitly labeled paired run otherwise.
\end{lstlisting}

\section{Representative Agent Output (excerpt)}
\begin{lstlisting}
thomas ncu: Duration 2.45ms, Compute 4%, Memory 55%,
Occupancy 9.95%, Waves/SM 0.18 -> GPU under-filled but each
thread is busy (M=1024 serial sweep), memory already coalesced.
Attempt 1 (PCR, 1 block/system): 20 barriers + 48KB smem ->
0.55x (SLOWER). Reverted.
Final: register-carry CSE (denominator once, carry U/RHS in
registers). CPU-verified bit-exact for M in {2..1024}.
Result: baseline 1.180ms / optimized 0.975ms -> 1.21x. PASS.
\end{lstlisting}

\section{Changes Made to the Agent}
The agent's guidance was iterated three times. v1 (free-form ``optimize this'')
produced confident but slower kernels (thomas PCR, gaussian fusion, jaccard
binary search). v2 added the mandatory \emph{profile-first / cite-a-metric} rule,
which stopped guess-driven changes. v3 added two rules that fixed the remaining
regressions: (i) \emph{prefer removing waste over adding work; never increase
total work to gain parallelism}, and (ii) \emph{a speedup may be reported only
after correctness passes}. We also added a same-compiler comparison harness and a
dedicated-node Slurm submission, after discovering that login-node timings and
mismatched \texttt{nvcc} versions both produced misleading speedups. Full
per-benchmark records (profile, attempts, code, correctness) are in the project
repository under \texttt{math\_data\_compress/profiling/} and
\texttt{optimized/}.

\end{document}
